% =====================================================================
\section{Conclusion}\label{sec:conclusion}
% =====================================================================

We have constructed a single $q^{2}$-dependent response kernel
$\kappa(q^{2})$ from a finite 2I-equivariant graph (the 600-cell
$V_{600}$, 120 vertices, 720 edges) regularised by the golden-ratio
mass $\varphi^{-2}$ as the discrete Green's response of the standard
graph Laplacian to a uniform equatorial-shell source, and tested
whether the same fixed shape — with one fitted dimensionless amplitude
per dataset — describes the $B\to K^{*}\mu^{+}\mu^{-}$ angular anomaly
as published by two collaborations across five fits and three decay
channels. The headline analysis evaluates every fit using
\texttt{flavio.np\_prediction} directly (no Taylor expansion in
$\Delta C_{9}$); the earlier linearised analysis is reported alongside
as a methodology diagnostic, with the difference between the two
modes the largest single uncertainty in the project.

The headline numerical results are:
\begin{itemize}
\item \textbf{LHCb 2025, four-observable joint fit (non-linear):}
  $A=+1.14$, $\Delta C_{9}^{\mathrm{eff}}=-0.86$, with
  $\Delta\mathrm{AIC}=+1.09$ relative to a constant-$\Delta C_{9}$
  shift at the same parameter cost. The kernel does not beat the
  constant shift on this dataset.
\item \textbf{Cross-dataset (non-linear, 5 fits across 2
  collaborations and 3 channels):} $\Delta\mathrm{AIC}\in
  [-0.24, +1.09]$ around zero, with $2/5$ favouring the kernel and
  $3/5$ favouring the constant shift. Stacked Akaike weight:
  $w_{\mathrm{VFD}}=0.35$, $w_{\mathrm{FREE\_C9}}=0.65$.
\item \textbf{Sign uniformity:} $A>0$ in all $5/5$ fits;
  $\Delta C_{9}^{\mathrm{eff}}<0$ in all $5/5$ fits. Under a random
  sign null this has $p=2^{-5}\approx 0.03$.
\item \textbf{Cross-dataset amplitude scatter:} non-linear
  $A\in[+1.05,+2.06]$ across the four $P$-basis $B\to K^{*}$ fits
  (factor $\sim 2$); $A=+4.98$ on the $S$-basis $B_{s}\!\to\!\phi$
  fit, mostly explained by the $P/S$ basis amplification factor
  $\langle 1/\sqrt{F_{L}(1-F_{L})}\rangle\approx 2.2$, with a
  factor-of-$2$ unexplained residual.
\item \textbf{Linearisation drift:} on LHCb 2025 the headline
  $\Delta\mathrm{AIC}$ shifts by $+2.77$ between the linearised and
  non-linear fits, larger than the linearised preference itself; the
  non-linear refit is taken as headline.
\item \textbf{Theoretical compression:} the kernel involves no fitted
  continuous shape parameters. The fit-tunable quantity per dataset
  is the dimensionless amplitude $A$, plus one
  pure-geometry-confirmed discrete variant choice
  (\S\ref{sec:derivation}, Table~\ref{tab:variant_selection}).
\end{itemize}

\paragraph{Falsification programme.} Each of the following events would
falsify the kernel-direction-consistency hypothesis:
\begin{enumerate}
\item A future $B\to K^{*}\mu\mu$ measurement (with any collaboration,
  any luminosity, any energy) gives $A<0$ on the geometry-derived kernel
  under non-linear refit.
\item Cross-dataset amplitude scatter on the $P$-basis fits exceeds
  a factor of $\sim 5$ once basis is harmonised.
\item Belle II angular $B\to K^{*}$ data, when released, gives an
  amplitude inconsistent with $+1.5\pm 0.5$ in the $P$-basis
  under non-linear refit.
\item A cross-channel $B^{+}\to K^{+}\mu\mu$ analysis (when public
  HEPData becomes available) gives $A<0$ under non-linear refit.
\end{enumerate}

The earlier project iteration's falsification trigger ``A future
analysis with non-linear $\Delta C_{9}$ evaluation gives
$\Delta\mathrm{AIC}_{\mathrm{VFD\,vs\,FREE\_C9}}\geq 0$ on the LHCb
2025 joint fit'' has now fired (\S\ref{subsec:joint_fit}). We
therefore explicitly weaken the model-comparison claim: the kernel is
sign-consistent with the anomaly but is not statistically preferred
over a constant Wilson-coefficient shift.

\paragraph{What this is, and is not.} This is a structural test of the
$b\to s\mu\mu$ angular anomaly's $q^{2}$ shape. \emph{What survives
the headline non-linear analysis:} a single geometry-derived kernel
with one fitted amplitude per dataset is consistent in direction with
the angular anomaly across two collaborations, two isospin partners,
and one cross-channel test. \emph{What does not survive:} an AIC
preference for the geometric shape over a constant
Wilson-coefficient shift. The data does not yet distinguish the
geometric kernel from a generic centre-peaked
nuisance.

\paragraph{Reproducibility.} All code, data, processed outputs, and
this paper are released under the MIT licence at
\citep{vfdBAnomalyRepo}. The complete pipeline (HEPData download,
SM-table regeneration via \texttt{flavio}, all driver scripts
including the non-linear refits and Akaike-weight stacking, figure
regeneration, paper recompilation) is reproducible from
\texttt{repro/run\_all.sh} on a clean checkout. Persistent
\texttt{flavio} caches under \texttt{data/processed/flavio\_cache.json}
ensure subsequent runs do not re-query \texttt{flavio} for already-computed
SM and non-linear-NP values. The non-linear refit scripts
(\texttt{scripts/wo016c\_nonlinear\_refit.py},
\texttt{scripts/wo016d\_nonlinear\_xdataset.py}) and the
Akaike-stacking script
(\texttt{scripts/wo016a\_akaike\_stack.py}) are stand-alone and
regenerate the headline tables of \S\ref{sec:results} and
\S\ref{sec:cross-dataset}.
